%%%%%%%%%%%%%%%%
%%% exod 21 %%%
%%%%%%%%%%%%%%%%

\Chapter{21}

\Verse{1}
{
	\hE{וְאֵלֶּה הַמִּשְׁפָּטִים אֲשֶׁר תָּשִׂים לִפְנֵיהֶם׃}
}
{
	\eE{
		“And these are the judgments that you shall set before\\
		them:\\
	}
}
{
%	\footnote{
%		these are the judgments that you shall set. These three chapters, Exodus
%		21–23, are a corpus of law known as the Covenant Code.
%	}
}

\Verse{2}
{
	\hE{כִּי תִקְנֶה עֶבֶד עִבְרִי שֵׁשׁ שָׁנִים יַעֲבֹד וּבַשְּׁבִעִת יֵצֵא לַחפְשִׁי חִנָּם׃}
}
{
	\eE{
		“When you will buy a Hebrew slave, he shall work six\\
		years, and in the seventh he shall go out liberated for\\
		free.\\
	}
}
{
%	\footnote{
%		Hebrew slave. In biblical narrative the term “Hebrew” is used to
%		identify Israelites only when speaking among foreigners. It is not the
%		standard term for the people, which is rather “Israelite” at first, and
%		“Jew” later. Its use here in biblical law, where we would have expected
%		“Israelite slave,” probably occurs because this had become a fixed
%		phrase through assonance: the two words, “slave” (, ‘ebed) and “Hebrew”
%		(, ‘ibrî), begin with the same two letters.
%	}
}

\Verse{3}
{
	\hE{אִם בְּגַפּוֹ יָבֹא בְּגַפּוֹ יֵצֵא אִם בַּעַל אִשָּׁה הוּא וְיָצְאָה אִשְׁתּוֹ עִמּוֹ׃}
}
{
	\eE{
		If he will come by himself, he shall go out by himself. If\\
		he is a woman’s husband, then his wife shall go out with\\
		him.\\
	}
}
{
}

\Verse{4}
{
	\hE{אִם אֲדֹנָיו יִתֶּן לוֹ אִשָּׁה וְיָלְדָה לוֹ בָנִים אוֹ בָנוֹת הָאִשָּׁה וִילָדֶיהָ תִּהְיֶה לַאדֹנֶיהָ וְהוּא יֵצֵא בְגַפּוֹ׃}
}
{
	\eE{
		If his master will give him a wife, and she will give\\
		birth to sons or daughters for him, the wife and her\\
		children will be her master’s, and he shall go out by\\
		himself.\\
	}
}
{
}

\Verse{5}
{
	\hE{וְאִם אָמֹר יֹאמַר הָעֶבֶד אָהַבְתִּי אֶת אֲדֹנִי אֶת אִשְׁתִּי וְאֶת בָּנָי לֹא אֵצֵא חפְשִׁי׃}
}
{
	\eE{
		And if the slave will say: ‘I love my master, my wife, and\\
		my children; I won’t go out liberated!’\\
	}
}
{
}

\Verse{6}
{
	\hE{וְהִגִּישׁוֹ אֲדֹנָיו אֶל הָאֱלֹהִים וְהִגִּישׁוֹ אֶל הַדֶּלֶת אוֹ אֶל הַמְּזוּזָה וְרָצַע אֲדֹנָיו אֶת אזְנוֹ בַּמַּרְצֵעַ וַעֲבָדוֹ לְעֹלָם׃}
}
{
	\eE{
		then his master shall bring him over to God and bring him\\
		over to the door or to the doorpost, and his master shall\\
		pierce his ear with an awl, and he shall serve him\\
		forever.\\
	}
}
{
}

\Verse{7}
{
	\hE{וְכִי יִמְכֹּר אִישׁ אֶת בִּתּוֹ לְאָמָה לֹא תֵצֵא כְּצֵאת הָעֲבָדִים׃}
}
{
	\eE{
		“And if a man will sell his daughter as a maid, she shall\\
		not go out as the slaves go out.\\
	}
}
{
}

\Verse{8}
{
	\hE{אִם רָעָה בְּעֵינֵי אֲדֹנֶיהָ אֲשֶׁר [לוֹ] (לא) יְעָדָהּ וְהֶפְדָּהּ לְעַם נכְרִי לֹא יִמְשֹׁל לְמכְרָהּ בְּבִגְדוֹ בָהּ׃}
}
{
	\eE{
		If she is bad in the eyes of her master who has designated\\
		her for himself, then he shall let her be redeemed. He\\
		shall not dominate so as to sell her to a foreign people\\
		in his betrayal of her.\\
	}
}
{
%	\footnote{
%		He shall not dominate. Recall that part of the curse on woman in Eden
%		was that “he’ll dominate you” (Gen 3:16). But here we are told of a case
%		in which a man is explicitly not free to dominate a woman. Even in that
%		worldview that recognized male control of women, there was an
%		understanding that this control is not unlimited. This point of respect
%		for women and recognition of a legal right that they possessed was a
%		first small step in the breakdown of that worldview. And this passage
%		teaches that a commandment can change or become limited even within the
%		Torah itself.
%	}
}

\Verse{9}
{
	\hE{וְאִם לִבְנוֹ יִיעָדֶנָּה כְּמִשְׁפַּט הַבָּנוֹת יַעֲשֶׂה לָּהּ׃}
}
{
	\eE{
		And if he will designate her for his son, he shall treat\\
		her according to the manner of daughters.\\
	}
}
{
}

\Verse{10}
{
	\hE{אִם אַחֶרֶת יִקַּח לוֹ שְׁאֵרָהּ כְּסוּתָהּ וְעֹנָתָהּ לֹא יִגְרָע׃}
}
{
	\eE{
		If he will take another for himself, he shall not subtract\\
		from her food, her apparel, and her hygiene.\\
	}
}
{
%	\footnote{
%		hygiene. Shalom Paul has shown parallels from ancient Near Eastern texts
%		that indicate that, beside providing food and clothing, the husband must
%		provide the necessary cosmetics for his wife’s cleanliness and
%		appearance. A second possibility is that the word here, ‘nth, is cognate
%		to Hebrew m‘ôn, meaning a dwelling, in which case the text means that
%		the husband provides food, clothing, and shelter. A third possibility is
%		the traditional view that this refers to sex, but there is no evidence
%		in support of this.
%	}
}

\Verse{11}
{
	\hE{וְאִם שְׁלשׁ אֵלֶּה לֹא יַעֲשֶׂה לָהּ וְיָצְאָה חִנָּם אֵין כָּסֶף׃}
}
{
	\eE{
		And if he will not do these three for her, then she shall\\
		go out free. There is no money.\\
	}
}
{
}

\Verse{12}
{
	\hE{מַכֵּה אִישׁ וָמֵת מוֹת יוּמָת׃}
}
{
	\eE{
		“One who strikes a man, and he dies, he shall be put to\\
		death.\\
	}
}
{
}

\Verse{13}
{
	\hE{וַאֲשֶׁר לֹא צָדָה וְהָאֱלֹהִים אִנָּה לְיָדוֹ וְשַׂמְתִּי לְךָ מָקוֹם אֲשֶׁר יָנוּס שָׁמָּה׃}
}
{
	\eE{
		And one who did not scheme, but God conveyed it to his\\
		hand, I shall set a place for you, that he shall flee\\
		there.\\
	}
}
{
}

\Verse{14}
{
	\hE{וְכִי יָזִד אִישׁ עַל רֵעֵהוּ לְהרְגוֹ בְערְמָה מֵעִם מִזְבְּחִי תִּקָּחֶנּוּ לָמוּת׃}
}
{
	\eE{
		But if a man will plot against his neighbor, to kill him\\
		with treachery, you shall take him from my altar to die.\\
	}
}
{
%	\footnote{
%		plot. The verb zyd elsewhere refers to cooking (e.g., Jacob’s cooking
%		the stew in Gen 25:29). Here it refers to someone’s committing murder
%		with intent, rather than manslaughter. (It also refers to the Egyptians’
%		scheme to enslave the Israelites in Exod 18:11.) Like the English
%		expression “to cook up,” the Hebrew word has the range of both cooking
%		and plotting.
%	}
}

\Verse{15}
{
	\hE{וּמַכֵּה אָבִיו וְאִמּוֹ מוֹת יוּמָת׃}
}
{
	\eE{
		“And one who strikes his father and his mother shall be\\
		put to death.\\
	}
}
{
}

\Verse{16}
{
	\hE{וְגֹנֵב אִישׁ וּמְכָרוֹ וְנִמְצָא בְיָדוֹ מוֹת יוּמָת׃}
}
{
	\eE{
		“And one who steals a man and has sold him, or he was\\
		found in his hand, will be put to death.\\
	}
}
{
}

\Verse{17}
{
	\hE{וּמְקַלֵּל אָבִיו וְאִמּוֹ מוֹת יוּמָת׃}
}
{
	\eE{
		“And one who curses his father and his mother shall be put\\
		to death.\\
	}
}
{
}

\Verse{18}
{
	\hE{וְכִי יְרִיבֻן אֲנָשִׁים וְהִכָּה אִישׁ אֶת רֵעֵהוּ בְּאֶבֶן אוֹ בְאֶגְרֹף וְלֹא יָמוּת וְנָפַל לְמִשְׁכָּב׃}
}
{
	\eE{
		“And if people will quarrel, and a man strikes his\\
		neighbor with a stone or with a fist, and he does not die,\\
		and he falls to bed,\\
	}
}
{
}

\Verse{19}
{
	\hE{אִם יָקוּם וְהִתְהַלֵּךְ בַּחוּץ עַל מִשְׁעַנְתּוֹ וְנִקָּה הַמַּכֶּה רַק שִׁבְתּוֹ יִתֵּן וְרַפֹּא יְרַפֵּא׃}
}
{
	\eE{
		if he will get up and walk himself outside on his staff,\\
		then the one who struck will be innocent. Only he shall\\
		compensate for his staying home and shall have him healed.\\
	}
}
{
}

\Verse{20}
{
	\hE{וְכִי יַכֶּה אִישׁ אֶת עַבְדּוֹ אוֹ אֶת אֲמָתוֹ בַּשֵּׁבֶט וּמֵת תַּחַת יָדוֹ נָקֹם יִנָּקֵם׃}
}
{
	\eE{
		“And if a man will strike his slave or his maid with a\\
		rod, and he dies by his hand, he shall be avenged.\\
	}
}
{
}

\Verse{21}
{
	\hE{אַךְ אִם יוֹם אוֹ יוֹמַיִם יַעֲמֹד לֹא יֻקַּם כִּי כַסְפּוֹ הוּא׃}
}
{
	\eE{
		Just: if he will stand for a day or two days, he shall not\\
		be avenged, because he is his money.\\
	}
}
{
%	\footnote{
%		stand. If he will not die immediately but will live for a day or two
%		after being struck. Footnote 21:21. he is his money. The slave is the
%		master’s property. This indicates that the master is empowered to strike
%		his slave, but he is not free to be the immediate cause of the slave’s
%		death.
%	}
}

\Verse{22}
{
	\hE{וְכִי יִנָּצוּ אֲנָשִׁים וְנָגְפוּ אִשָּׁה הָרָה וְיָצְאוּ יְלָדֶיהָ וְלֹא יִהְיֶה אָסוֹן עָנוֹשׁ יֵעָנֵשׁ כַּאֲשֶׁר יָשִׁית עָלָיו בַּעַל הָאִשָּׁה וְנָתַן בִּפְלִלִים׃}
}
{
	\eE{
		“And if people will fight, and they strike a pregnant\\
		woman, and her children go out, and there will not be an\\
		injury, he shall be penalized according to what the\\
		woman’s husband will impose on him, and he will give it by\\
		the judges.\\
	}
}
{
%	\footnote{
%		her children go out, and there will not be an injury. This can mean: if
%		the woman has a miscarriage, but there is no other physical damage to
%		the woman herself. Or it can mean: if the woman goes into labor, and the
%		child is born without any physical damage. Scholars generally understand
%		it to be the former. This passage has been cited in debates over
%		abortion. However, the issues in this case are different, and the
%		wording is so complex, and its meaning is so uncertain, that one should
%		be extremely hesitant to use this case for either side of that debate.
%		Footnote 21:22. he will give it by the judges. The meaning is uncertain,
%		but the penalty is determined and enforced in some way involving both
%		the husband and the judges.
%	}
}

\Verse{23}
{
	\hE{וְאִם אָסוֹן יִהְיֶה וְנָתַתָּה נֶפֶשׁ תַּחַת נָפֶשׁ׃}
}
{
	\eE{
		And if there will be an injury, then you shall give a life\\
		for a life, 21:23–24. a life for a life, an eye for an eye\\
		…The inclusion of “a burn for a burn” here—in a case in\\
		which “burn” has no apparent relevance—indicates that this\\
		list comes to apply a fundamental principle to this case.\\
		Concerning the eye-for-an-eye principle, see the comment\\
		on Lev 24:20.\\
	}
}
{
}

\Verse{24}
{
	\hE{עַיִן תַּחַת עַיִן שֵׁן תַּחַת שֵׁן יָד תַּחַת יָד רֶגֶל תַּחַת רָגֶל׃}
}
{
	\eE{
		an eye for an eye, a tooth for a tooth, a hand for a hand,\\
		a foot for a foot,\\
	}
}
{
}

\Verse{25}
{
	\hE{כְּוִיָּה תַּחַת כְּוִיָּה פֶּצַע תַּחַת פָּצַע חַבּוּרָה תַּחַת חַבּוּרָה׃}
}
{
	\eE{a burn for a burn, a wound for a wound, a hurt for a hurt.}
}
{
}

\Verse{26}
{
	\hE{וְכִי יַכֶּה אִישׁ אֶת עֵין עַבְדּוֹ אוֹ אֶת עֵין אֲמָתוֹ וְשִׁחֲתָהּ לַחפְשִׁי יְשַׁלְּחֶנּוּ תַּחַת עֵינוֹ׃}
}
{
	\eE{
		“And if a man will strike his slave’s eye or his maid’s\\
		eye and destroy it, he shall let him go, liberated, for\\
		his eye.\\
	}
}
{
%	\footnote{
%		liberated, for his eye. Even in a system that allows slavery, a master
%		does not own the slave’s body. The master can require work from a slave,
%		but damaging the slave’s eye or tooth steps over the line, and the
%		master no longer owns the slave. This reminds us also of the Egyptian
%		enslavement of Israel. The Egyptians are not criticized for having
%		slaves, but rather for the way they treat the slaves: “They degraded
%		them…. They made them serve with harshness … they made their lives
%		bitter with hard work …” They killed the male infants. The Torah did not
%		bring about the end of slavery by abolishing it. It established
%		principles regarding slaves’ dignity, rights, and treatment. And these
%		gradually contributed to humans’ own rejection of slavery in much
%		(though not yet all) of the world. This is a crucial point itself: that
%		some of the laws of the Torah command things outright, while others lead
%		humans to grow and change themselves.
%	}
}

\Verse{27}
{
	\hE{וְאִם שֵׁן עַבְדּוֹ אוֹ שֵׁן אֲמָתוֹ יַפִּיל לַחפְשִׁי יְשַׁלְּחֶנּוּ תַּחַת שִׁנּוֹ׃}
}
{
	\eE{
		And if he will knock out his slave’s tooth or his maid’s\\
		tooth, he shall let him go liberated for his tooth.\\
	}
}
{
}

\Verse{28}
{
	\hE{וְכִי יִגַּח שׁוֹר אֶת אִישׁ אוֹ אֶת אִשָּׁה וָמֵת סָקוֹל יִסָּקֵל הַשּׁוֹר וְלֹא יֵאָכֵל אֶת בְּשָׂרוֹ וּבַעַל הַשּׁוֹר נָקִי׃}
}
{
	\eE{
		“And if an ox will gore a man or a woman and they die, the\\
		ox shall be stoned, and its meat shall not be eaten—and\\
		the ox’s owner is innocent.\\
	}
}
{
}

\Verse{29}
{
	\hE{וְאִם שׁוֹר נַגָּח הוּא מִתְּמֹל שִׁלְשֹׁם וְהוּעַד בִּבְעָלָיו וְלֹא יִשְׁמְרֶנּוּ וְהֵמִית אִישׁ אוֹ אִשָּׁה הַשּׁוֹר יִסָּקֵל וְגַם בְּעָלָיו יוּמָת׃}
}
{
	\eE{
		And if it was a goring ox from the day before yesterday,\\
		and it had been so testified to its owner, and he did not\\
		watch it, and it killed a man or a woman, the ox will be\\
		stoned, and its owner will be put to death as well.\\
	}
}
{
%	\footnote{
%		from the day before yesterday. Meaning: this ox had gored before.
%	}
}

\Verse{30}
{
	\hE{אִם כֹּפֶר יוּשַׁת עָלָיו וְנָתַן פִּדְיֹן נַפְשׁוֹ כְּכֹל אֲשֶׁר יוּשַׁת עָלָיו׃}
}
{
	\eE{
		If a ransom will be set on him, then he shall give\\
		everything that will be set on him for the redemption of\\
		his life.\\
	}
}
{
}

\Verse{31}
{
	\hE{אוֹ בֵן יִגָּח אוֹ בַת יִגָּח כַּמִּשְׁפָּט הַזֶּה יֵעָשֶׂה לּוֹ׃}
}
{
	\eE{
		Whether it will gore a son or it will gore a daughter,\\
		according to this judgment shall be done to it. 21:31–32.\\
		a son or … a daughter … a slave or a maid. No distinction\\
		is made on the basis of sex; a distinction is made on the\\
		basis of class: free versus slave.\\
	}
}
{
}

\Verse{32}
{
	\hE{אִם עֶבֶד יִגַּח הַשּׁוֹר אוֹ אָמָה כֶּסֶף שְׁלֹשִׁים שְׁקָלִים יִתֵּן לַאדֹנָיו וְהַשּׁוֹר יִסָּקֵל׃}
}
{
	\eE{
		If the ox will gore a slave or a maid, he shall pay thirty\\
		silver shekels to his owner, and the ox shall be stoned.\\
	}
}
{
}

\Verse{33}
{
	\hE{וְכִי יִפְתַּח אִישׁ בּוֹר אוֹ כִּי יִכְרֶה אִישׁ בֹּר וְלֹא יְכַסֶּנּוּ וְנָפַל שָׁמָּה שּׁוֹר אוֹ חֲמוֹר׃}
}
{
	\eE{
		“And if a man will open a pit, or if a man will dig a pit,\\
		and he will not cover it, and an ox or an ass will fall\\
		there,\\
	}
}
{
}

\Verse{34}
{
	\hE{בַּעַל הַבּוֹר יְשַׁלֵּם כֶּסֶף יָשִׁיב לִבְעָלָיו וְהַמֵּת יִהְיֶה לּוֹ׃}
}
{
	\eE{
		the owner of the pit shall pay; he shall pay back money to\\
		its owner, and the dead one will be his.\\
	}
}
{
}

\Verse{35}
{
	\hE{וְכִי יִגֹּף שׁוֹר אִישׁ אֶת שׁוֹר רֵעֵהוּ וָמֵת וּמָכְרוּ אֶת הַשּׁוֹר הַחַי וְחָצוּ אֶת כַּסְפּוֹ וְגַם אֶת הַמֵּת יֶחֱצוּן׃}
}
{
	\eE{
		And if a man’s ox will strike his neighbor’s ox, and it\\
		dies, then they shall sell the live ox and split the money\\
		for it, and they shall split the dead one as well.\\
	}
}
{
}

\Verse{36}
{
	\hE{אוֹ נוֹדַע כִּי שׁוֹר נַגָּח הוּא מִתְּמוֹל שִׁלְשֹׁם וְלֹא יִשְׁמְרֶנּוּ בְּעָלָיו שַׁלֵּם יְשַׁלֵּם שׁוֹר תַּחַת הַשּׁוֹר וְהַמֵּת יִהְיֶה לּוֹ׃}
}
{
	\eE{
		Or: if it was known that it was a goring ox from the day\\
		before yesterday, and its owner did not watch it, he shall\\
		pay an ox for the ox, and the dead one shall be his.\\
	}
}
{
}

\Verse{37}
{
	\hOther{כִּי יִגְנֹב אִישׁ שׁוֹר אוֹ שֶׂה וּטְבָחוֹ אוֹ מְכָרוֹ חֲמִשָּׁה בָקָר יְשַׁלֵּם תַּחַת הַשּׁוֹר וְאַרְבַּע צֹאן תַּחַת הַשֶּׂה׃}
}
{
	\eOther{
		“If a man will steal an ox or a sheep and slaughter it or\\
		sell it, he shall pay five oxen for the ox and four sheep\\
		for the sheep.\\
	}
}
{
}
